% GENERATED FILE. Edit canonical lessons, then run npm run generate:books.
% canonical-source-hash: fnv1a64:b950abd9a08544f1
% canonical-lessons: RU-C338-argument, RU-C338-dovod, RU-C338-dokazyvat, RU-C338-dokazatelstvo, RU-C338-ubeditelnyy

\chapter{Arguments and evidence}
\label{ch:ru-arguments-and-evidence}
\hlchaptermodality{338}

\begin{chapteropening}
\textbf{By the end of this chapter:} \emph{I can give an argument and a reason, prove a point with evidence and say whether it is convincing.}

Complete the last lesson of chapter 338: I can give an argument and a reason, prove a point with evidence and say whether it is convincing.
\end{chapteropening}

\section[\texorpdfstring{\ru{аргумент} (argumént)}{аргумент (argumént)}]{\ru{аргумент} (argumént) — an argument (a reason)}

\label{lesson:RU-C338-argument}



\begin{quote}
Before the new one: say the Russian for to declare, then the Russian for to be considered.
\end{quote}

\subsection*{You'll want to know: \ru{аргумент}}
\textbf{\ru{аргумент}} (\emph{argumént}) — \textquotedblleft{}an argument (a reason)\textquotedblright{}.

\begin{quote}
\ru{Это} \ru{сильный} \ru{аргумент}.
\end{quote}

\begin{quote}
\emph{Éto sílnyy argumént.}
\end{quote}

\begin{quote}
That is a strong argument.
\end{quote}

Masculine. A reason in a debate, not a quarrel: a quarrel is \emph{\ru{спор}} or \emph{\ru{ссора}}.

\begin{cousinweb}
\textbf{\ru{аргумент}} is usually traced, through German and Polish, to Latin \emph{argūmentum} \textquotedblleft{}proof, evidence\textquotedblright{} (Lewis \& Short), built on \emph{arguere} \textquotedblleft{}to make clear, to prove\textquotedblright{}.
\end{cousinweb}

\begin{grammarlens}[title={What you've built}]
One more word for telling, explaining and describing at B1.
\end{grammarlens}

\subsection*{Your turn}
\begin{itemize}
\raggedright
  \item \emph{Say it:} \emph{argumént}
  \item \emph{Say it:} \emph{argumént}, once more
  \item \emph{Say it:} a sentence of your own with \emph{argumént}
  \item \emph{Recall:} say the Russian for to declare, then the Russian for to be considered, then say \emph{argumént} again
\end{itemize}

\begin{tcolorbox}[breakable,colback=teal!4,colframe=teal!35!black,title={Before you move on}]
What does \ru{аргумент} mean? (\textquotedblleft{}An argument (a reason)\textquotedblright{}.) Say it once more.
\end{tcolorbox}
\section[\texorpdfstring{\ru{довод} (dóvod)}{довод (dóvod)}]{\ru{довод} (dóvod) — a point made, a reason given}

\label{lesson:RU-C338-dovod}



\begin{quote}
Before the new one: say the Russian for to be considered, then the Russian for an argument.
\end{quote}

\subsection*{You'll want to know: \ru{довод}}
\textbf{\ru{довод}} (\emph{dóvod}) — \textquotedblleft{}a point made, a reason given\textquotedblright{}.

\begin{quote}
\ru{Его} \ru{доводы} \ru{меня} \ru{не} \ru{убедили}.
\end{quote}

\begin{quote}
\emph{Yegó dóvody menyá ne ubedíli.}
\end{quote}

\begin{quote}
His points did not convince me.
\end{quote}

Masculine. More bookish than \emph{\ru{аргумент}}: \emph{\ru{привести} \ru{довод}} — to put forward a point.

\begin{cousinweb}
\textbf{\ru{довод}} is built from \emph{\ru{до}-} \textquotedblleft{}up to\textquotedblright{} and \emph{\ru{вод}-}, the root of \emph{\ru{водить}} \textquotedblleft{}to lead\textquotedblright{}, the iterative partner of \emph{\ru{вести}}, from Proto-Slavic \emph{*voditi} \textquotedblleft{}to lead, to conduct\textquotedblright{} (EtymDB). A point made leads the listener up to a conclusion.
\end{cousinweb}

\begin{grammarlens}[title={What you've built}]
One more word for telling, explaining and describing at B1.
\end{grammarlens}

\subsection*{Your turn}
\begin{itemize}
\raggedright
  \item \emph{Say it:} \emph{dóvod}
  \item \emph{Say it:} \emph{dóvod}, once more
  \item \emph{Say it:} a sentence of your own with \emph{dóvod}
  \item \emph{Recall:} say the Russian for to be considered, then the Russian for an argument, then say \emph{dóvod} again
\end{itemize}

\begin{tcolorbox}[breakable,colback=teal!4,colframe=teal!35!black,title={Before you move on}]
What does \ru{довод} mean? (\textquotedblleft{}A point made, a reason given\textquotedblright{}.) Say it once more.
\end{tcolorbox}
\section[\texorpdfstring{\ru{доказывать} (dokázyvat)}{доказывать (dokázyvat)}]{\ru{доказывать} (dokázyvat) — to prove}

\label{lesson:RU-C338-dokazyvat}



\begin{quote}
Before the new one: say the Russian for an argument, then the Russian for a point made.
\end{quote}

\subsection*{You'll want to know: \ru{доказывать}}
\textbf{\ru{доказывать}} (\emph{dokázyvat}) — \textquotedblleft{}to prove\textquotedblright{}.

\begin{quote}
\ru{Он} \ru{долго} \ru{доказывал}, \ru{что} \ru{прав}.
\end{quote}

\begin{quote}
\emph{On dólgo dokázyval, chto prav.}
\end{quote}

\begin{quote}
He spent a long time proving that he was right.
\end{quote}

Imperfective, with the perfective \emph{\ru{доказать}} (\emph{\ru{докажу}, \ru{докажешь}}). \emph{\ru{Доказывать}, \ru{что}…} — to prove that…

\begin{cousinweb}
\textbf{\ru{доказывать}} is built from \emph{\ru{до}-} and \emph{\ru{казать}} \textquotedblleft{}to show\textquotedblright{}, the verb inside \emph{\ru{показывать}} and \emph{\ru{сказать}}. \emph{\ru{Казать}} is inherited from Proto-Slavic \emph{*kazati} (EtymDB): to prove is to show a thing through to the end.
\end{cousinweb}

\begin{grammarlens}[title={What you've built}]
One more word for telling, explaining and describing at B1.
\end{grammarlens}

\subsection*{Your turn}
\begin{itemize}
\raggedright
  \item \emph{Say it:} \emph{dokázyvat}
  \item \emph{Say it:} \emph{dokázyvat}, once more
  \item \emph{Say it:} a sentence of your own with \emph{dokázyvat}
  \item \emph{Recall:} say the Russian for an argument, then the Russian for a point made, then say \emph{dokázyvat} again
\end{itemize}

\begin{tcolorbox}[breakable,colback=teal!4,colframe=teal!35!black,title={Before you move on}]
What does \ru{доказывать} mean? (\textquotedblleft{}To prove\textquotedblright{}.) Say it once more.
\end{tcolorbox}
\section[\texorpdfstring{\ru{доказательство} (dokazátelstvo)}{доказательство (dokazátelstvo)}]{\ru{доказательство} (dokazátelstvo) — proof, evidence}

\label{lesson:RU-C338-dokazatelstvo}



\begin{quote}
Before the new one: say the Russian for a point made, then the Russian for to prove.
\end{quote}

\subsection*{You'll want to know: \ru{доказательство}}
\textbf{\ru{доказательство}} (\emph{dokazátelstvo}) — \textquotedblleft{}proof, evidence\textquotedblright{}.

\begin{quote}
\ru{У} \ru{полиции} \ru{нет} \ru{доказательств}.
\end{quote}

\begin{quote}
\emph{U polítsii net dokazátelstv.}
\end{quote}

\begin{quote}
The police have no evidence.
\end{quote}

Neuter. The plural \emph{\ru{доказательства}} means pieces of evidence.

\begin{cousinweb}
\textbf{\ru{доказательство}} is built on \emph{\ru{доказать}} \textquotedblleft{}to prove\textquotedblright{} with the abstract suffix \emph{-\ru{тельство}}; the root is \emph{\ru{казать}} \textquotedblleft{}to show\textquotedblright{}, from Proto-Slavic \emph{*kazati} (EtymDB).
\end{cousinweb}

\begin{grammarlens}[title={What you've built}]
One more word for telling, explaining and describing at B1.
\end{grammarlens}

\subsection*{Your turn}
\begin{itemize}
\raggedright
  \item \emph{Say it:} \emph{dokazátelstvo}
  \item \emph{Say it:} \emph{dokazátelstvo}, once more
  \item \emph{Say it:} a sentence of your own with \emph{dokazátelstvo}
  \item \emph{Recall:} say the Russian for a point made, then the Russian for to prove, then say \emph{dokazátelstvo} again
\end{itemize}

\begin{tcolorbox}[breakable,colback=teal!4,colframe=teal!35!black,title={Before you move on}]
What does \ru{доказательство} mean? (\textquotedblleft{}Proof, evidence\textquotedblright{}.) Say it once more.
\end{tcolorbox}
\section[\texorpdfstring{\ru{убедительный} (ubedítelnyy)}{убедительный (ubedítelnyy)}]{\ru{убедительный} (ubedítelnyy) — convincing}

\label{lesson:RU-C338-ubeditelnyy}



\begin{quote}
Before the new one: say the Russian for to prove, then the Russian for proof.
\end{quote}

\subsection*{You'll want to know: \ru{убедительный}}
\textbf{\ru{убедительный}} (\emph{ubedítelnyy}) — \textquotedblleft{}convincing\textquotedblright{}.

\begin{quote}
\ru{Его} \ru{ответ} \ru{был} \ru{не} \ru{очень} \ru{убедительным}.
\end{quote}

\begin{quote}
\emph{Yegó otvét byl ne óchen ubedítelnym.}
\end{quote}

\begin{quote}
His answer was not very convincing.
\end{quote}

From \emph{\ru{убедить}}, to convince. \emph{\ru{Убедительная} \ru{просьба}} — an urgent request, a fixed formula on notices.

\begin{cousinweb}
\textbf{\ru{убедительный}} is built on \emph{\ru{убедить}} \textquotedblleft{}to convince\textquotedblright{}. The root is that of \emph{\ru{беда}} \textquotedblleft{}trouble\textquotedblright{}, inherited from Proto-Slavic \emph{*běda}, which EtymDB sets beside Greek \emph{peíthō} \textquotedblleft{}I persuade\textquotedblright{} (EtymDB). To convince someone was first to press them hard.
\end{cousinweb}

\begin{grammarlens}[title={What you've built}]
One more word for telling, explaining and describing at B1.
\end{grammarlens}

\subsection*{Your turn}
\begin{itemize}
\raggedright
  \item \emph{Say it:} \emph{ubedítelnyy}
  \item \emph{Say it:} \emph{ubedítelnyy}, once more
  \item \emph{Say it:} a sentence of your own with \emph{ubedítelnyy}
  \item \emph{Recall:} say the Russian for to prove, then the Russian for proof, then say \emph{ubedítelnyy} again
\end{itemize}

\begin{tcolorbox}[breakable,colback=teal!4,colframe=teal!35!black,title={Before you move on}]
What does \ru{убедительный} mean? (\textquotedblleft{}Convincing\textquotedblright{}.) Say it once more.
\end{tcolorbox}
